\section{Background} \label{sec:background}

\ap{provide a illustation for CRDTs}
\subsection{Sequential Types}
We begin the discussion with a definition of a sequential type.
A sequential type captures the behavior of an object whose state mutates by applying operations one at a time. 
Each operation transforms the current state into the next, starting from a well defined initial state and following a deterministic transition relation that governs an object's behavior. 
The sequence of operations on a sequential type---called a history---is a total order.
Formally, we define a sequential type as follows: 

\begin{definition}[Sequential Type]
A sequential type $T = (S,s_0, \Delta)$  is a triple
    where $S$ is the set of states of an object of type $T$, $s_0 \in S $ the initial state and $\Delta: S \times Op \to S$ is the transition function mapping a state and an operation to a successor state. 
    We write $s \xrightarrow{op_1 \dots op_n} s'$ when $\Delta(s, \{op_1 \dots op_n\}) = s'$. 
\end{definition}

We can define a binary operator ($\oplus$) over the set of states that a sequential type can achieve.
This binary operator allows folding two states into one while preserving the intention of operations used to achieve said states. 
Formally, we define this binary operator as follows:

\begin{definition}[Intention-preserving]
  Operator $\oplus$ is said \emph{intention-preserving} when for any two states $s$ and $s'$, for any two sequences of operations $\sigma$ and $\sigma'$ such that $\Delta(s_0,\sigma)=s$ and $\Delta(s_0,\sigma')=s'$, there exists a sequences $\lambda$ with $ops(\lambda)=ops(\sigma) \cup ops(\sigma')$ such that $\Delta(s_0,\lambda)=s \oplus s'$
\end{definition}

\begin{definition}[Induced binary operator]
  Given a sequential type $T$, a binary operator $\oplus: S \times S \to S$ is defined such that for states $s , s', s'' \in S$, $s \oplus s' = s''$ such that $s''$ preserves the effects of operations that achieved $s$ and $s'$ from $s_0$. 
If $s_0 \xrightarrow{op_1 \dots  op_m} s$ and $s_0 \xrightarrow{op_1' \dots  op_n'} s'$. Then, $s \oplus s' = \Delta (s_0, \{op_1 \dots op_m\} \uplus \{op_1'\dots op_n'\})$.
\end{definition}

\ap{discuss the condition for well-definedness via confluence.}

\sd{For $\oplus$ to be well defined, the result must not depend on the particular
linearization of the merged operations. We therefore assume throughout that the
transition system is confluent, i.e., whenever two sequences
of operations are reorderings of one another, they reach the same state. Under
confluence, $\Delta^*$ is independent of the order in which operations are
applied, and $\oplus$ is a total function on $S$.}



\subsection{Concurrent Types and Eventual Consistency}
Unlike a sequential type, that restricts operations to execute one after the other a \emph{concurrent type} allows multiple operations to take effect simultaneously. 
This sequence of operations is a partial order. 
Since concurrent types admit partially ordered histories, they can be used to type shared objects. 

Shared objects operate under a consistency model. 
To byapss the limitations of CAP theorem and support availability, Eventual consistency is often the preferred model. 
Eventual consistency states that, \sd{once all updates have been
delivered,} all parties apply the same set of operations on a shared object.
Eventual consistency is strengthened with guaranteed convergence and termination to define \emph{Strong eventual consistency} \ap{cite paper}. 

A shared type that satisfies strong eventual consistency is called a \emph{Conflict-free Replicated Datatype} (CRDT). 
A CRDT is defined via a state representation, a set of query operations to read the state, a set of update operations, and a merge operator to merge any two states.
A replica independently updates its local copy of a CRDT.
In order to include remote updates, it merges the remote states into the local copy.
Convergence is guaranteed by requiring the merge operation to be associative, commutative, and idempotent.
Formally, a CRDT is defined as follows.

\begin{definition}[CRDT]
    A conflict-free replicated datatype, denoted $\mathbb{T} \coloneq (\mathbb{S}, \le \bot, u, \sqcup)$, is a tuple consisting of a set of states $\mathbb{S}$, containing an initial state $\bot \in \mathbb{S}$. 
    Other states are derived from this initial state via a set of update operations $u$. 
    CRDT state is monotonically non-decreasing across updates i.e. $u(s)=s' \Rightarrow s \le s'$.     
    A binary merge operator $\sqcup : \mathbb{S} \times \mathbb{S} \rightarrow \mathbb{S}$ computes the lowest upper bound of two states in $\mathbb{S}$. 
    Merge induces a partial ordering over the states where $s_1 \sqcup s_2 = s_3 \Rightarrow s_1 \le s_3$ and $s_2 \le s_3$

    % \ps{In this definition $\leq$ is the ordering implied by $\sqcup$ (see PG-CRDT paper). Thus $\leq$ can be omitted.}
\end{definition}



% Formally, a state-based CRDT is defined by a tuple $(S,O,Q,R,\tau,\Merge)$ where %
% $S$ is a set of states;
% $O$ is a set of update operations;
% $R$ is a set of responses;
% $Q$ is a set of query operations ($S \times Q \rightarrow R$);
% $\tau$ is a total transition function $S \times O \rightarrow S$; and
% $\Merge \in S \times S \rightarrow S$ is a merge operator.
% %
% It is required that $(S,\Merge)$ forms a join-semilattice, to ensure convergence.
% Moreover, the state ordering induced by $\Merge$ is compatible with the transition function, i.e., $x \leq y$ whenever $x \Merge y = y$ and $x \leq \tau(x,o)$, for any update $o$.



\begin{example}[Counter CRDTs] \label{ex:counterCRDT}
  A basic CRDT example is a grow-only counter.
 % For each replica $r$, its internal state maps to an integer $C(r)$ that tracks the number of increments performed by $r$.
  Each replica $r$ maintains an integer $C(r)$ counting its increments.
  Whenever $r$ invokes operation $incr()$, it increments its entry $C(r)$ by one.
  Reading the counter returns the sum of all entries in $C$.
  Merging states $C$ and $C'$ takes the maximum entry per replica: $C \Merge C'(r) = \mathit{max}\{C(r), C'(r)\}$.
  
  A composed CRDT is the \emph{Positive-Negative Counter}; PN-counter for short.
  Internally, it is implemented as a pair of two grow-only counters $(P, N)$.
  It supports $incr()$ and $decr()$, storing increments in $P$ and decrements in $N$.
  The current value of a PN-counter reads as $P - N$.
  Merging computes $(P \Merge P', N \Merge N')$.
  
  To illustrate, suppose a system with two replicas $r$ and $r'$.
   Replica $r$ calls $incr()$ five times and $decr()$ once, reaching the internal state $(P: \{r: 5, r': 0\}, N: \{r: 1, r': 0\})$.
  Concurrently, $r'$ calls $decr()$ twice, reaching $(P: \{r: 0, r': 0\}, N:\{r: 0, r': 2\})$.
  Merging the two states yields $(P:\{r: 5, r': 0\}, N:\{r: 1, r': 2\})$, which represents the value $2$.
\end{example}


\subsection{Correctness of CRDTs}
\ap{Present the current discussion on the approaches used to define the correctness of CRDTs. }

\sd{Two complementary notions of correctness recur in the literature. The first
is convergence, i.e., any two replicas that have observed the same set of
updates read the same value, which SEC guarantees by requiring
$(\mathbb{S}, \sqcup)$ to form a join-semilattice~\cite{DBLP:conf/sss/ShapiroPBZ11}.
The second is behavioral adequacy, i.e., the replicated object should realize the
intended sequential specification, observed through a materialization (or
observation) function~\cite{DBLP:journals/ftpl/Burckhardt14,DBLP:conf/popl/BurckhardtGYZ14}. Our
notion of liftability, developed next, makes this second requirement precise by
tying the observation function to the induced operator $\oplus$ of the sequential
type.}
